\SourceLeaf{18}
woraus $r=\BareRoot{((a-x)^2+(b-y)^2+(c-z)^2)}$ folgt
\[
\begin{aligned}
\cos MX&=\frac{a-x}{r}\\
\cos MY&=\frac{b-y}{r}\\
\cos MZ&=\frac{c-z}{r}
\end{aligned}
\]
Schließlich gilt nach einem wohlbekannten Lehrsatz
\[
\cos MQ=\cos MX.\cos QX+\cos MY.\cos QY+\cos MZ.\cos QZ
\]
oder
\[
\cos MQ=\pm\frac{T(a-x)+U(b-y)+V(c-z)}{r\BareRoot{(TT+UU+VV)}}
\]

\Article{10}
\par
Um die Differentialausdrücke über die ganze Oberfläche integrieren zu können, müssen wir sie nun so umformen, dass sie nur zwei Variablen enthalten. Dies ist zwar möglich, indem man eine der Variablen $x,y,z$ mithilfe der Gleichung $W=0$ eliminiert, doch sind die so gewonnenen Formeln meist schwieriger zu handhaben. Daher ist es besser, zwei neue Unbestimmte $p,q$ einzuführen und sowohl $x$ als auch $y$ und $z$ als Funktionen dieser Unbestimmten anzusehen.

\begingroup\rightskip=0pt plus 1em
Sobald man also $p,q$ bestimmte Werte gibt, sind auch $x,y,z$ bestimmt. Diesen Werten entspricht also ein bestimmter Punkt auf der Oberfläche des Körpers. Diese wechselseitige Zuordnung wird deutlicher, wenn wir uns eine unbegrenzte Ebene denken, deren Punkte durch rechtwinklige Koordinaten $p,q$ gegeben sind. Jedem Punkt der Ebene entspricht dann ein Punkt auf der Oberfläche des Körpers, und zwar genau einer, wenn die Zuordnung so eingerichtet ist, dass $x,y,z$ eindeutige Funktionen der Unbestimmten $p,q$ sind. Wenn umgekehrt durch $x,y,z$ auch $p$ und $q$ vollständig und eindeutig bestimmt sind, entspricht offenbar jedem Punkt der Körperoberfläche genau ein Punkt der Ebene. In diesem Fall muss sich die Ebene nach allen Seiten ins Unendliche erstrecken, damit sie die gesamte Körperoberfläche erfasst. Andernfalls brauchen wir nur einen Teil der Ebene zu betrachten, der durch endliche oder unendliche Grenzen bestimmt ist und die Oberfläche des Körpers gewissermaßen darstellt. Man denke sich die Ebene durch unendlich viele Geraden, die teils parallel zur Abszissenachse und teils senkrecht zu ihr verlaufen,
